% Section VII (part) -- Experimental Analysis: predictive results (Person 2)
% Include inside Section VII of main.tex with:  % Section VII (part) -- Experimental Analysis: predictive results (Person 2)
% Include inside Section VII of main.tex with:  % Section VII (part) -- Experimental Analysis: predictive results (Person 2)
% Include inside Section VII of main.tex with:  % Section VII (part) -- Experimental Analysis: predictive results (Person 2)
% Include inside Section VII of main.tex with:  \input{sections/predictive_results}
%
% The table and figure are GENERATED -- do not type numbers by hand:
%   python -m experiments.compare_models --data fi2010 --seeds 42 43 44
%   python -m experiments.confusion_matrices --seed 42
% then copy results/classification_results.tex and figures/confusion_all.pdf
% next to main.tex (or adjust the paths below).
%
% Every sentence marked TODO must be rewritten from the real FI-2010 numbers.
\subsection{Predictive Performance (Experiments 1 and 3)}
\label{sec:predictive}

Table~\ref{tab:classification} reports test-set accuracy, macro-F1 and
macro-precision for the three models at horizon $k=10$, averaged over
three seeds (mean $\pm$ standard deviation). Because the Stationary class
dominates FI-2010 at short horizons, a classifier that always predicts
Stationary already attains high accuracy; we therefore treat macro-F1 as the
primary metric.

\input{results/classification_results}

\textbf{Effect of adding attention (CNN vs.\ CNN + standard attention).}
% TODO: state the macro-F1 difference and whether it exceeds the seed std.
Adding one softmax-attention encoder block changes macro-F1 by
\textbf{TODO} points relative to the CNN alone, at the cost of
33{,}472 additional parameters.

\textbf{Effect of the attention mechanism (standard vs.\ linear).}
% TODO: compare rows 2 and 3. "Competitive" = within ~1 std across seeds.
Replacing softmax attention with linear attention changes macro-F1 by
\textbf{TODO} points. Combined with the efficiency results in
Section~\ref{sec:efficiency}, this quantifies the accuracy--efficiency
trade-off of the proposed model.

\begin{figure*}[t]
  \centering
  \includegraphics[width=0.8\textwidth]{figures/confusion_all.pdf}
  \caption{Row-normalised test-set confusion matrices (seed 42, $k=10$).
  Diagonal entries are per-class recall; raw counts in parentheses.}
  \label{fig:confusion}
\end{figure*}

\textbf{Error analysis.} Figure~\ref{fig:confusion} shows where each model
errs.
% TODO: describe from results/per_class_report.txt, e.g. whether Up and Down
% are confused with each other (direction errors) or mostly with Stationary
% (magnitude errors), and which model recovers the minority classes best.

%
% The table and figure are GENERATED -- do not type numbers by hand:
%   python -m experiments.compare_models --data fi2010 --seeds 42 43 44
%   python -m experiments.confusion_matrices --seed 42
% then copy results/classification_results.tex and figures/confusion_all.pdf
% next to main.tex (or adjust the paths below).
%
% Every sentence marked TODO must be rewritten from the real FI-2010 numbers.
\subsection{Predictive Performance (Experiments 1 and 3)}
\label{sec:predictive}

Table~\ref{tab:classification} reports test-set accuracy, macro-F1 and
macro-precision for the three models at horizon $k=10$, averaged over
three seeds (mean $\pm$ standard deviation). Because the Stationary class
dominates FI-2010 at short horizons, a classifier that always predicts
Stationary already attains high accuracy; we therefore treat macro-F1 as the
primary metric.

\input{results/classification_results}

\textbf{Effect of adding attention (CNN vs.\ CNN + standard attention).}
% TODO: state the macro-F1 difference and whether it exceeds the seed std.
Adding one softmax-attention encoder block changes macro-F1 by
\textbf{TODO} points relative to the CNN alone, at the cost of
33{,}472 additional parameters.

\textbf{Effect of the attention mechanism (standard vs.\ linear).}
% TODO: compare rows 2 and 3. "Competitive" = within ~1 std across seeds.
Replacing softmax attention with linear attention changes macro-F1 by
\textbf{TODO} points. Combined with the efficiency results in
Section~\ref{sec:efficiency}, this quantifies the accuracy--efficiency
trade-off of the proposed model.

\begin{figure*}[t]
  \centering
  \includegraphics[width=0.8\textwidth]{figures/confusion_all.pdf}
  \caption{Row-normalised test-set confusion matrices (seed 42, $k=10$).
  Diagonal entries are per-class recall; raw counts in parentheses.}
  \label{fig:confusion}
\end{figure*}

\textbf{Error analysis.} Figure~\ref{fig:confusion} shows where each model
errs.
% TODO: describe from results/per_class_report.txt, e.g. whether Up and Down
% are confused with each other (direction errors) or mostly with Stationary
% (magnitude errors), and which model recovers the minority classes best.

%
% The table and figure are GENERATED -- do not type numbers by hand:
%   python -m experiments.compare_models --data fi2010 --seeds 42 43 44
%   python -m experiments.confusion_matrices --seed 42
% then copy results/classification_results.tex and figures/confusion_all.pdf
% next to main.tex (or adjust the paths below).
%
% Every sentence marked TODO must be rewritten from the real FI-2010 numbers.
\subsection{Predictive Performance (Experiments 1 and 3)}
\label{sec:predictive}

Table~\ref{tab:classification} reports test-set accuracy, macro-F1 and
macro-precision for the three models at horizon $k=10$, averaged over
three seeds (mean $\pm$ standard deviation). Because the Stationary class
dominates FI-2010 at short horizons, a classifier that always predicts
Stationary already attains high accuracy; we therefore treat macro-F1 as the
primary metric.

\input{results/classification_results}

\textbf{Effect of adding attention (CNN vs.\ CNN + standard attention).}
% TODO: state the macro-F1 difference and whether it exceeds the seed std.
Adding one softmax-attention encoder block changes macro-F1 by
\textbf{TODO} points relative to the CNN alone, at the cost of
33{,}472 additional parameters.

\textbf{Effect of the attention mechanism (standard vs.\ linear).}
% TODO: compare rows 2 and 3. "Competitive" = within ~1 std across seeds.
Replacing softmax attention with linear attention changes macro-F1 by
\textbf{TODO} points. Combined with the efficiency results in
Section~\ref{sec:efficiency}, this quantifies the accuracy--efficiency
trade-off of the proposed model.

\begin{figure*}[t]
  \centering
  \includegraphics[width=0.8\textwidth]{figures/confusion_all.pdf}
  \caption{Row-normalised test-set confusion matrices (seed 42, $k=10$).
  Diagonal entries are per-class recall; raw counts in parentheses.}
  \label{fig:confusion}
\end{figure*}

\textbf{Error analysis.} Figure~\ref{fig:confusion} shows where each model
errs.
% TODO: describe from results/per_class_report.txt, e.g. whether Up and Down
% are confused with each other (direction errors) or mostly with Stationary
% (magnitude errors), and which model recovers the minority classes best.

%
% The table and figure are GENERATED -- do not type numbers by hand:
%   python -m experiments.compare_models --data fi2010 --seeds 42 43 44
%   python -m experiments.confusion_matrices --seed 42
% then copy results/classification_results.tex and figures/confusion_all.pdf
% next to main.tex (or adjust the paths below).
%
% Every sentence marked TODO must be rewritten from the real FI-2010 numbers.
\subsection{Predictive Performance (Experiments 1 and 3)}
\label{sec:predictive}

Table~\ref{tab:classification} reports test-set accuracy, macro-F1 and
macro-precision for the three models at horizon $k=10$, averaged over
three seeds (mean $\pm$ standard deviation). Because the Stationary class
dominates FI-2010 at short horizons, a classifier that always predicts
Stationary already attains high accuracy; we therefore treat macro-F1 as the
primary metric.

\input{results/classification_results}

\textbf{Effect of adding attention (CNN vs.\ CNN + standard attention).}
% TODO: state the macro-F1 difference and whether it exceeds the seed std.
Adding one softmax-attention encoder block changes macro-F1 by
\textbf{TODO} points relative to the CNN alone, at the cost of
33{,}472 additional parameters.

\textbf{Effect of the attention mechanism (standard vs.\ linear).}
% TODO: compare rows 2 and 3. "Competitive" = within ~1 std across seeds.
Replacing softmax attention with linear attention changes macro-F1 by
\textbf{TODO} points. Combined with the efficiency results in
Section~\ref{sec:efficiency}, this quantifies the accuracy--efficiency
trade-off of the proposed model.

\begin{figure*}[t]
  \centering
  \includegraphics[width=0.8\textwidth]{figures/confusion_all.pdf}
  \caption{Row-normalised test-set confusion matrices (seed 42, $k=10$).
  Diagonal entries are per-class recall; raw counts in parentheses.}
  \label{fig:confusion}
\end{figure*}

\textbf{Error analysis.} Figure~\ref{fig:confusion} shows where each model
errs.
% TODO: describe from results/per_class_report.txt, e.g. whether Up and Down
% are confused with each other (direction errors) or mostly with Stationary
% (magnitude errors), and which model recovers the minority classes best.
